% Part: second-order-logic
% Chapter: syntax-and-semantics
% Section: language-of-sol

\documentclass[../../../include/open-logic-section]{subfiles}

\begin{document}

\olfileid{sol}{syn}{lan}

\olsection{ద్వితీయ క్రమ తర్కం యొక్క భాష}

ప్రథమ క్రమ తర్కంలో మాదిరిగానే, ద్వితీయ క్రమ తర్కంలోని వ్యక్తీకరణలను \emph{\tetoken{చరాలు}{!!{variable}s}}, \emph{\tetoken{స్థిరాంకాలు}{!!{constant}s}}, \emph{\tetoken{విధేయాలు}{!!{predicate}s}}, కొన్నిసార్లు \emph{\tetoken{ప్రమేయాలు}{!!{function}s}} ఉన్న ప్రాథమిక పదసంపద నుంచి నిర్మిస్తాం. వీటికి తార్కిక సంయోజకాలు, పరిమాణకాలు, కుండలీకరణాలు, కామాలు వంటి విరామ చిహ్నాలను జోడించి \emph{పదాలను}, \emph{\tetoken{సూత్రాలను}{!!{formula}s}} రూపొందిస్తాం. తేడా ఏమిటంటే, వస్తువులకు సంబంధించిన చరాలతోపాటు ద్వితీయ క్రమ తర్కంలో సంబంధాలకు, ప్రమేయాలకు సంబంధించిన చరాలు కూడా ఉంటాయి; వాటిపై పరిమాణీకరణకు అనుమతి ఉంటుంది. కాబట్టి ద్వితీయ క్రమ తర్కంలోని తార్కిక చిహ్నాలు ప్రథమ క్రమ తర్కంలోని చిహ్నాలతోపాటు కిందివి:

\begin{enumerate}
\item ప్రతి స్థానసంఖ్య~$n$కు ద్వితీయ క్రమ సంబంధ \tetoken{చరాల}{!!{variable}s} \tetoken{గణనీయ అనంత}{!!{denumerable}s} సమితి: $\Obj V_0^n$, $\Obj V_1^n$, $\Obj V_2^n$, \dots
\item ప్రతి స్థానసంఖ్యకు ద్వితీయ క్రమ ప్రమేయ \tetoken{చరాల}{!!{variable}s} \tetoken{గణనీయ అనంత}{!!{denumerable}s} సమితి: $\Obj u_0^n$, $\Obj u_1^n$, $\Obj u_2^n$, \dots
\end{enumerate}

ప్రథమ క్రమ తర్కంలో సాధారణంగా \tetoken{తాదాత్మ్య చిహ్నం}{!!{identity}}~$\eq$ కూడా ఉంటుంది. ప్రథమ క్రమ తర్కంలో ఆసక్తికరమైన విషయాలను వ్యక్తపరచడానికి భాష~$\Lang{L}$ యొక్క తర్కేతర చిహ్నాలు ముఖ్యమైనవి. తర్కేతర చిహ్నాలను ఉపయోగించని \tetoken{వాక్యాలు}{!!{sentences}} కూడా ఉంటాయి; కానీ కేవలం~$\eq$తో ఆసక్తికరమైన విషయాలు చెప్పడం కష్టం. ద్వితీయ క్రమ తర్కంలో సంబంధ, ప్రమేయ చరాలు అపరిమితంగా లభిస్తాయి కాబట్టి, ప్రథమ క్రమ భాషలోని తర్కేతర చిహ్నాలకు బదులు తగిన స్వేచ్ఛా సంబంధ, ప్రమేయ చరాలను ఉపయోగించి వాటికి ఆ చిహ్నాల వ్యాఖ్యానాలను విలువలుగా కేటాయిస్తే, ప్రత్యేక తర్కేతర చిహ్నాల సరఫరా లేకుండానే అదే ప్రతిపాదనలను వ్యక్తపరచవచ్చు.\sourcecorrection{OLTESOLSYNLAN-001}{మూలంలోని చివరి వాదనకు అవసరమైన విలువ కేటాయింపు నిబంధనను స్పష్టం చేశాం. తర్కేతర చిహ్నాల స్థానంలో స్వేచ్ఛా చరాలను ఉపయోగించడం, వాటిని పరిమాణకంతో బంధించి మూసిన వాక్యాన్ని ఏర్పరచడం ఒకటి కాదు; మూల చిహ్నాల వ్యాఖ్యానాలను సంబంధిత చరాలకు కేటాయించినప్పుడే ఈ ప్రతిస్థాపన అదే విషయాన్ని వ్యక్తపరుస్తుంది.}

\end{document}
